% 参苓白术散.tex

\documentclass[12pt,UTF8]{ctexbook}

% 设置纸张信息。
\input{../../../Book/common/page}

% 设置字体，并解决显示难检字问题。
\xeCJKsetup{AutoFallBack=true}
% 注意：ctexbook类已默认设置SimSun为CJKrmdefault
% 以下设置用于确保字体回退和扩展字体可用
% 仅设置扩展字体以避免与默认设置冲突
\setCJKfamilyfont{hei}{SimHei}
\setCJKfamilyfont{kai}{KaiTi}

% 目录 chapter 级别加点（.）。
\usepackage{titletoc}
\titlecontents{chapter}[0pt]{\vspace{3mm}\bf\addvspace{2pt}\filright}{\contentspush{\thecontentslabel\hspace{0.8em}}}{}{\titlerule*[8pt]{.}\contentspage}

% 支持目录点击跳转
\usepackage[colorlinks,linkcolor=blue,citecolor=blue,urlcolor=blue]{hyperref}

% 设置 part 和 chapter 标题格式。
\ctexset{
	part/name= {第,卷},
	part/number={\chinese{part}},
	chapter/name={第,篇},
	chapter/number={\arabic{chapter}}
}

% 图片相关设置。
\usepackage{graphicx}
\graphicspath{{Images/}}

% 设置署名格式。
\newenvironment{shuming}{\hfill\zihao{4}}

% 注脚每页重新编号，避免编号过大。
\usepackage[perpage]{footmisc}

% 设置古文原文格式。
\newenvironment{yuanwen}{\bfseries\zihao{4}}

% 设置段落间距
\setlength{\parskip}{0.8em plus 0.2em minus 0.1em}

% 列表格式支持，列表项向右偏移。
\usepackage{enumitem}

% 定义方剂组成环境
\newenvironment{fangji}[1]{\section{#1}\begin{enumerate}[leftmargin=2cm]}{\end{enumerate}}

\title{\heiti\zihao{0} 参苓白术散}
\author{WangFei}
\date{\today}

\begin{document}

\maketitle
\tableofcontents

\frontmatter
\chapter{前言}

\mainmatter

\chapter{参苓白术散}

\section{基本信息}
\begin{itemize}
	\item 【别名】：参术饮、白术散
	\item 【出处】：《太平惠民和剂局方》
	\item 【分类】：补益剂-补气
\end{itemize}

\section{组成}
\begin{fangji}{药物组成}
	\item 人参 10g
	\item 茯苓 15g
	\item 白术（麸炒）15g
	\item 山药 15g
	\item 白扁豆（炒）12g
	\item 莲子 10g
	\item 薏苡仁（炒）15g
	\item 砂仁 6g
	\item 桔梗 10g
	\item 甘草 6g
\end{fangji}

\section{性状}
本品为淡黄色的水丸；气香，味甜。

\section{用法}
\begin{itemize}
	\item 【煎服方法】：水煎服，每日1剂，分2次温服
	\item 【服用剂量】：成人用量10-15g，儿童酌减
	\item 【服用时间】：宜饭前服用或进食同时服
	\item 【特殊处理】：无特殊处理
\end{itemize}

\section{功效}
\begin{itemize}
	\item 【主要功效】：补脾胃，益肺气
	\item 【作用机制】：通过益气健脾、渗湿止泻，改善脾胃功能
\end{itemize}

\section{主治}
\begin{itemize}
	\item 【传统主治】：脾胃虚弱，食少便溏，气短咳嗽，肢倦乏力
	\item 【现代主治】：慢性胃肠炎、贫血、慢性支气管炎、慢性肾炎、妇女带下病等属脾虚湿盛者
\end{itemize}

\section{方解}
\begin{description}
	\item[君药]：人参、白术、茯苓 - 益气健脾渗湿
	\item[臣药]：山药、莲子 - 健脾益气，兼能止泻
	\item[佐药]：白扁豆、薏苡仁 - 健脾渗湿；砂仁 - 醒脾和胃，行气化滞；桔梗 - 宣肺利气，载药上行
	\item[使药]：甘草 - 健脾和中，调和诸药
\end{description}

\section{临床应用}
\begin{itemize}
	\item 【辨证要点】：食少便溏，肢倦乏力，舌淡苔白腻，脉虚缓
	\item 【现代应用】：慢性胃肠炎、消化不良、肠易激综合征、慢性支气管炎、慢性肾炎等
	\item 【病例举隅】：适用于长期脾胃虚弱，反复腹泻，神疲乏力的患者
\end{itemize}

\section{加减变化}
\begin{fangji}{随证加减}
	\item 兼气滞者，加陈皮、木香
	\item 兼寒者，加干姜、肉桂
	\item 兼热者，加黄连、黄芩
	\item 兼湿重者，加苍术、藿香
\end{fangji}

\section{注意事项}
\begin{fangji}{使用禁忌}
	\item 泄泻兼有大便不通畅，肛门有下坠感者忌服
	\item 服本药时不宜同时服用藜芦、五灵脂、皂荚或其制剂
	\item 不宜喝茶和吃萝卜以免影响药效
	\item 不宜和感冒类药同时服用
	\item 高血压、心脏病、肾脏病、糖尿病严重患者及孕妇应在医师指导下服用
\end{fangji}

\section{附方}
\begin{itemize}
	\item 【相关方剂】：四君子汤、六君子汤、香砂六君子汤
	\item 【比较分析】：四君子汤为基础方，参苓白术散在此基础上加山药、莲子、白扁豆等增强健脾渗湿作用
\end{itemize}

\backmatter

\end{document}